%!TEX root = ../../main.tex
\section{선행 연구 대비 기여의 성격}
\label{sec:disc-contribution}

선행 연구 \cite{baek2026killconditioned}와 비교했을 때 본 논문의 기여는 ``더 높은 승자 AUC''도 ``새로운 심층 아키텍처''도 아니다. 기여는 (1) 동결 승률 평가기 아래의 경기 연결 교전 후 가치 변화 정의, (2) 측정 품질($V \to W$)과 방향 예측 가능성($q \to \SVI$)을 분리하는 검증 스택, (3) 이득이 유연한 $(\ppre, t)$ 기준선, 작은 변화 필터, 이후 환경에서 살아남는지의 명시적 시험이다. 결과는 ``범위가 한정된 예측 가능성과 전이 한계''로 요약되며, 강력하고 이식 가능한 교전 예측 엔진의 주장이 아니다. 라벨과 코호트가 달라졌으므로 선행 연구의 AUC와 본 논문의 AUC 사이의 차이를 개선으로 부르지 않는다.

학위논문으로서의 깊이는 같은 증거를 설계 근거와 대안, 그리고 한계 중심으로 전개한 데 있다. 각 상수의 출처 분류(제 \ref{sec:eng-constants} 절), 규모 절단에 따른 부호 반전(제 \ref{sec:eng-scale} 절), 킬 없는 전투의 측정 불가능성(제 \ref{sec:eng-killless} 절), 두 평가기 경로(제 \ref{sec:data-roles} 절), 실행이 설계와 달랐던 부분의 기록(제 \ref{sec:pred-baselines} 절)은 저널 원고가 압축한 부분이다.
